```latex
\documentclass[11pt]{article}

\usepackage{amsmath,amsthm,amssymb,mathtools}
\usepackage[colorlinks=true,linkcolor=blue,citecolor=blue,urlcolor=blue]{hyperref}

\newtheorem{theorem}{Theorem}
\newtheorem{lemma}[theorem]{Lemma}
\newtheorem{remark}[theorem]{Remark}

\newcommand{\Q}{\mathbb{Q}}
\newcommand{\Z}{\mathbb{Z}}
\newcommand{\HF}{\widehat{HF}}

\title{Surgeries on L-space knots have at most two prime summands}
\author{}
\date{}

\begin{document}

\maketitle

\begin{abstract}
We prove that if $K\subset S^3$ is an L-space knot and $r$ is any Dehn surgery slope, then the surgery manifold $S^3_r(K)$ has at most two prime summands. The proof is a short deduction from Hoffman's theorem on reducible surgeries on L-space knots, together with standard facts about $0$-surgery on fibered knots and the Gordon--Luecke integrality theorem.
\end{abstract}

\paragraph{Convention.}
Throughout, an \emph{L-space knot} means a knot in $S^3$ admitting a positive Dehn surgery to an L-space. We write $S^3_r(K)$ for the closed $3$-manifold obtained by Dehn surgery on $K$ with slope $r\in \Q\cup\{\infty\}$.

\begin{theorem}[Hoffman {\cite{Hoffman}}]\label{thm:Hoffman}
Let $K\subset S^3$ be an L-space knot. If $S^3_r(K)$ is reducible, then $S^3_r(K)$ has at most two irreducible summands.
\end{theorem}

The theorem just cited is the essential Heegaard--Floer-theoretic input. To convert it into the desired statement about \emph{prime} summands, we only need to rule out the possibility of an $S^1\times S^2$ factor. For L-space knots this is straightforward.

\begin{lemma}\label{lem:zero-irred}
If $K\subset S^3$ is a nontrivial L-space knot, then $S^3_0(K)$ is irreducible.
\end{lemma}

\begin{proof}
By Ozsv\'ath--Szab\'o's description of knot Floer homology for knots admitting L-space surgeries \cite{OS} and Ni's fibred-knot detection theorem \cite{Ni}, every nontrivial L-space knot is fibered.

Let $F$ be a fiber surface for $K$. Then the knot exterior $S^3\setminus \nu K$ is the mapping torus of a homeomorphism $\varphi\colon F\to F$. Performing $0$-surgery caps off $\partial F$ by a disk, so $S^3_0(K)$ is the mapping torus of the induced homeomorphism of the closed surface
\[
\widehat F = F\cup_{\partial F} D^2.
\]
Since $K$ is nontrivial, $g(\widehat F)=g(K)\ge 1$. A surface bundle over $S^1$ with closed fiber of genus at least $1$ is irreducible; see, for example, Hempel \cite[Chapter~10]{Hempel}. Hence $S^3_0(K)$ is irreducible.
\end{proof}

We can now prove the desired statement.

\begin{theorem}\label{thm:main}
Let $K\subset S^3$ be an L-space knot, and let $r\in \Q\cup\{\infty\}$. Then $S^3_r(K)$ has at most two prime summands.
\end{theorem}

\begin{proof}
If $r=\infty$, then $S^3_r(K)=S^3$, which is prime.

If $K$ is the unknot, then every surgery on $K$ is a lens space (or $S^1\times S^2$ in the case $r=0$), so again there is only one prime summand.

Hence assume $K$ is nontrivial and $r\in \Q$. If $S^3_r(K)$ is irreducible, there is nothing to prove. So suppose that $S^3_r(K)$ is reducible.

By Lemma~\ref{lem:zero-irred}, the slope $r$ cannot be $0$. On the other hand, Gordon and Luecke proved that a reducible surgery on a knot in $S^3$ must have \emph{integral} slope \cite{GL}. Therefore
\[
r\in \Z\setminus\{0\}.
\]
It follows that
\[
H_1\!\left(S^3_r(K);\Z\right)\cong \Z/|r|\Z
\]
is finite.

Now recall the standard fact from prime decomposition theory that the only orientable prime closed $3$-manifold which is not irreducible is $S^1\times S^2$; see, e.g., Hempel \cite{Hempel}. Since
\[
H_1(S^1\times S^2;\Z)\cong \Z
\]
is infinite, $S^1\times S^2$ cannot occur as a prime summand of $S^3_r(K)$. Consequently, for this manifold the prime summands and irreducible summands are the same.

Applying Hoffman's theorem (Theorem~\ref{thm:Hoffman}), we conclude that $S^3_r(K)$ has at most two irreducible summands, hence at most two prime summands.
\end{proof}

\begin{remark}
If one adopts the broader convention that also calls the mirror of a positive L-space knot an ``L-space knot,'' the same conclusion still holds, because
\[
S^3_r(\overline K)\cong -\,S^3_{-r}(K),
\]
and orientation reversal does not change the number of prime summands.
\end{remark}

\begin{remark}
For context, the torus and satellite cases are already classical. Moser proved that the only reducible surgery on a nontrivial torus knot $T_{p,q}$ is $pq$-surgery, and
\[
S^3_{pq}(T_{p,q})\cong L(p,q)\# L(q,p)
\]
\cite{Moser}. Scharlemann proved that if a satellite knot admits a reducible surgery, then it is a cable knot and the reducing slope is the cabling slope \cite{Scharlemann}; Gordon's cabling formula identifies that surgery as
\[
L(p,q)\# S^3_{q/p}(J)
\]
for the companion knot $J$ \cite{GordonSatellite}, and Gordon--Luecke imply that the second factor is irreducible because $q/p\notin \Z$ \cite{GL}. Thus the genuinely nonclassical input needed for Theorem~\ref{thm:main} is precisely Hoffman's Heegaard--Floer-theoretic result.
\end{remark}

\begin{thebibliography}{99}

\bibitem{GL}
C.~McA. Gordon and J.~Luecke,
\emph{Only integral Dehn surgeries can yield reducible manifolds},
Math. Proc. Cambridge Philos. Soc. \textbf{102} (1987), no.~1, 97--101.

\bibitem{GordonSatellite}
C.~McA. Gordon,
\emph{Dehn surgery and satellite knots},
Trans. Amer. Math. Soc. \textbf{275} (1983), no.~2, 687--708.

\bibitem{Hempel}
J.~Hempel,
\emph{3-Manifolds},
Annals of Mathematics Studies, No.~86,
Princeton University Press, Princeton, N.J., 1976.

\bibitem{Hoffman}
C.~M. Hoffman,
\emph{Reducible surgeries and Heegaard Floer homology},
Algebr. Geom. Topol. \textbf{16} (2016), no.~6, 3431--3456.

\bibitem{Moser}
L.~Moser,
\emph{Elementary surgery along a torus knot},
Pacific J. Math. \textbf{38} (1971), 737--745.

\bibitem{Ni}
Y.~Ni,
\emph{Knot Floer homology detects fibred knots},
Invent. Math. \textbf{170} (2007), no.~3, 577--608.

\bibitem{OS}
P.~Ozsv\'ath and Z.~Szab\'o,
\emph{On knot Floer homology and lens space surgeries},
Topology \textbf{44} (2005), no.~6, 1281--1300.

\bibitem{Scharlemann}
M.~Scharlemann,
\emph{Producing reducible $3$-manifolds by surgery on a knot},
Topology \textbf{29} (1990), no.~4, 481--500.

\end{thebibliography}

\end{document}
```